% 本文件为论文主 TeX 文件
% 使用 texlive 完整编译:
% xelatex -> bibtex -> xelatex -> xelatex
% UNIV-Thesis 通用的本科/硕博毕业论文 LaTeX 模板

\documentclass[master]{univthesis}
% bachelor 本科毕业论文, 不能省略
% master 硕士学位论文, 默认可省略
% doctor 博士学位论文, 不能省略
% bachelor 本科毕业论文, 不能省略
% print 用于打印, 封面等生成空白页

% fonts macOS系统必选, 调用本地字体
% 需在本地文件夹 fonts 中添加字体 SimSun.ttf, KaiTi.ttf, SimHei.ttf, FangSong.ttf

% 这是一个通用的本科/硕/博毕业论文 LaTeX 模板
% 可以通过更换封面与声明页的 PDF 文件适配不同学校
% 这里封面和声明页可以通过 Word 转为 PDF 文件插入
% 可以修改类文件以满足不同格式要求

%----- 设置英文字体 -----
%\usepackage{newtxtext}  % New TX font for text
%\setmainfont{TeX Gyre Termes}  % Times New Roman 的开源复刻版本
%\setsansfont{TeX Gyre Heros}   % Helvetica 的开源复刻版本
%\setmonofont{TeX Gyre Cursor}  % Courier New 的开源复刻版本
\setmainfont{Times New Roman}
\setsansfont{Arial}
%\setmonofont{Courier New}

%----- 设置数学字体 -----
%\usepackage{newtxmath}
%\usepackage{mathptmx}

%----- 添加其他宏包 -----
%\usepackage[notcite,notref]{showkeys}
\usepackage{listings}
\usepackage{subfig}

%----- 设置超链接颜色 -----
%\hypersetup{linkcolor=magenta}
%\hypersetup{hidelinks}

%----- 参考文献格式 -----
%\bibliographystyle{plain} % abbrv, unsrt, siam
\bibliographystyle{univthesis-numeric}
%\bibliographystyle{univthesis-author-year}

%----- 参考文献引用格式 -----
\citestyle{numbers}
%\citestyle{super}
%\citestyle{authoryear}

%----- 调整列表项的间距 -----
%\setlength{\itemsep}{3pt}

%----- 调整页面避免出现过大空白 -----
%\raggedbottom

%----- 定义符号描述命令  -----
\newcommand{\nameditem}[3][]{
\noindent\hspace{2em}\makebox[0.2\textwidth][l]{#2}{{#3}
\hfill\makebox[0.2\textwidth][l]{#1}\hspace*{2em}}\par}

%----- 微分算子 -----
\newcommand*{\dif}{\mathop{}\!\mathrm{d}}

%----- 自定义命令 -----
\newcommand{\CC}{\ensuremath{\mathbb{C}}}
\newcommand{\RR}{\ensuremath{\mathbb{R}}}
\newcommand{\bA}{\boldsymbol{A}}
\newcommand{\ii}{\mathrm{i}\mkern1mu}
\newcommand{\abs}[1]{\lvert#1\rvert}
\newcommand{\norm}[1]{\left\lVert#1\right\rVert}
\newcommand{\dx}[1][x]{\mathop{}\!\mathrm{d}#1}
\newcommand{\red}[1]{\textcolor{red}{#1}}

\begin{document}


% 插入中英文封面, 可用 Word 生成
% 中文封面
\insertcover{cover-zh.pdf}
% 英文封面
\insertcover{cover-en.pdf}

% 插入独创性与授权声明页, 可用 Word 生成
\insertstatement{statement.pdf}

% 以上两个命令都使用了 pdfpages 宏包的 includepdf 命令插入 PDF 文件
% \includepdf[选项]{PDF文件名}


%=============== 摘要 =================%

% 关闭章节编号, 页码使用罗马数字
\frontmatter

% 中文摘要内容和关键字
\input{data/cnabstract}

% 英文摘要内容和关键字
\input{data/enabstract}

%=============== 目录 =================%

% 生成目录
% bachelor 选项下, 目录本身不会加入目录页
\maketoc

% 生成插图清单, 如不需要可以注释
% bachelor 选项下, 插图清单不会加入目录页
\makelof

% 生成表格清单, 如不需要可以注释
% bachelor 选项下, 表格清单不会加入目录页
\makelot

%=====================================%

% 主要符号表, 如不需要可以注释
\input{data/denotation}

%============== 正文部分 ==============%

% 开启章节编号, 页码使用阿拉伯数字
\mainmatter

% 第1章 引言
\input{data/chap01}

% 第2章 LaTeX 常用环境
\input{data/chap02}

% 第3章 微分方程的数值方法
\input{data/chap03}

% 第4章 插图环境
\input{data/chap04}

% 第5章 表格环境
\input{data/chap05}


%============== 参考文献 ==============%

% 生成参考文献, 两种方式任选一种

% 第一种方式, 使用 bib 文件
%\nocite{*}  % 可以显示全部参考文献
\bibliography{reference}

%------------------------------------%

% 第二种方式, 手动添加文献信息
%\input{data/bibliography}

%=============== 附录 ================%

% 添加附录, 如不需要可以注释
\input{data/appendix}


%====================================%

\backmatter % 结束章节自动编号

% 本科/攻读学位期间的研究成果, 如不需要可以注释
\input{data/publications}

% 致谢
\input{data/acknowledgement}


\end{document}
